\scanpage{81}
donc $a_0$ serait fixée par $\sigma_0$ et $\sigma_1$, donc par $G$
donc par les $u^n$, donc $\varphi_a$ serait constante, cqfd.
En $=$ temps, on a retrouvé (pour mémoire)
(9), qui montre que les données (2) (3)
satisfaisant (8) $+$ conditions lemme équivalent à
(2), \deleted{\unsure{($3_s$) (savoir|($a_s$) (savoir}} \inserted{$+$ donnée} $s_0$) satisfaisant que $s_0,s_1$ et $s_{-1}$

\uline{Proposition} Les réalisations géom régulières de $\Pi_\infty$
dans $E$ correspondent 1-1 aux triples
$(\sigma_0,\sigma_1,s_0)$, où
\[ \left\{\begin{aligned} &\sigma_0,\sigma_1\in\operatorname{Aff}(E)\\
&s_0\in E\end{aligned}\right. \tag*{10)} \]
satisfaisant les conditions

a) $\sigma_0^2=\sigma_1^2=\operatorname{id}$

b) $\sigma_1s_0=s_0$ i.e. $s_0\in E^{\sigma_1}$

c) $s_0,\sigma_0s_0=s_1$, et $s_{-1}\overset{\unsure{\mathrm{dfn}|\mathrm{déf}|\mathrm{def}}}{=}\sigma_1s_1$ forment une
base affine de $E$.

La repr. $\varphi_G:G=\operatorname{Aut}(\Pi_{\unsure{+|\infty}})\simeq\unsure{Ga_1|G_{a1}|G_1|G_{0,1}}$ associée est
définie par
\[ \varphi_G(\underline{\sigma}_0)=\sigma_0,\quad \varphi_G(\underline{\sigma}_1)=\sigma_1 \tag*{(\unsure{11|10})} \]
et $\varphi_s$ et $\varphi_a$ sont donnés par
\[ \left.\begin{aligned}\varphi_s(\underline s_n)&=s_n\overset{\mathrm{dfn}}{=}u^n s_0\\
\varphi_a(\underline a_n)&=\operatorname{dr}(s_n,s_{n+1})=u^n\operatorname{dr}(s_0,s_1)\end{aligned}\right|\, n\in\mathbb Z \tag*{\(\begin{aligned}&(12)\\ &(\unsure{13|12})\end{aligned}\)} \]
où
\[ u=\sigma_0\sigma_1 \tag*{(14)} \]
% archive sheet number: 80
